The same compensated native source admits a joint energy and charge-arrival
measurement. Let $E=P_I(H)$ be its complete interacting energy bin and let
$F_e$ resolve a native arrival route by the original local charged pole.
The extra excited link gives $H_0F_e=(a+1)F_e$. Consequently
$\|F_eE\|\leq v/g$, where $v=\|\lambda T\|$ and $g$ is the distance
from $a+1$ to the bin. The actual charged resolvents fix the selected
spin-one residue tensor and yield the first-hop moment
$m=w^3(1/12+x_{\rm cube}/648)>0$. A commutator identity bounds
$F_eEBP_0$ by $-\lambda F_eTS_aBP_0$ with a complete $O(\lambda^2)$
remainder. This requires no global isolated charged-band hypothesis.
The physical preparation error enters as $15\lambda\eta$ in amplitude,
so the source's small nonzero probability remains resolvable.
Each route has original-input time-zero probability strictly between
$1.04\times10^{-23}$ and $1.69\times10^{-23}$.

The tagged joint probability is positive at $t=0$ after energy
filtering, whereas $F_eBP_0=0$ without that filter. It is therefore a
virtual tail resolved by measurement, not established propagation. The
same Hamiltonian bounds its subsequent change relative to this baseline.
Seventeen complete records retain preparation, pump and energy failures,
tagged and untagged arrivals, and the full nonarrival complement. The
atomic tag can change the measured state's energy; it is not assumed
to commute with the interacting energy bin. No autonomous free-electron
mass shell or continuum particle identification follows.
